% 35-29(35쪽) — 원점을 지나 증가하며 위로 볼록한 함수 y=f(x) 의 그래프와 직선 y=2x. 스캔 화질이 낮아 TikZ 로 다시 그림.
% 원문에 있는 것만 옮긴다: 곡선 · 직선 · a·b 에서 곡선까지 올라가는 안내 점선 둘 · 라벨 y=2x, y=f(x), O, a, b, x, y.
\begin{tikzpicture}[x=0.95cm, y=0.95cm, line width=0.6pt]
  \draw[->, >=stealth, line width=0.7pt] (-0.95,0) -- (4.45,0) node[below=1pt, font=\small] {$x$};
  \draw[->, >=stealth, line width=0.7pt] (0,-1.25) -- (0,3.05) node[left=1pt, font=\small] {$y$};
  \draw[densely dotted, line width=0.4pt] (1.90,0) -- (1.90,1.68);
  \draw[densely dotted, line width=0.4pt] (2.95,0) -- (2.95,2.11);
  \draw[line width=0.6pt] (-0.55,-1.10) -- (1.35,2.70);
  \draw[line width=0.6pt, domain=0:4.20, samples=120, smooth] plot (\x, {2.75*(1-exp(-0.50*\x))});
  \node[anchor=south west, font=\small] at (1.05,2.55) {$y=2x$};
  \node[anchor=south, font=\small] at (3.35,2.35) {$y=f(x)$};
  \node[below right=-1pt and -1pt, font=\small] at (0,0) {$\pt{O}$};
  \node[below=1pt, font=\small] at (1.90,0) {$a$};
  \node[below=1pt, font=\small] at (2.95,0) {$b$};
\end{tikzpicture}
